% part: intuitionistic-logic
% chapter: soundness-completeness
% section: decidability

\documentclass[../../../include/open-logic-section]{subfiles}

\begin{document}

\olfileid{int}{sc}{dec}

\olsection{القابلية للقرار}

كلما كان~$\Gamma \Proves/ !A$ أعطانا برهان التمام نموذجًا عوالمه غير
متناهية، يشهد بـ~$\Gamma \Entails/ !A$. أما النتيجة الآتية فتبيّن أن
إقامة~$\Entails !A$ لا تحتاج إلا إلى إقامة
$\mSat{M}{!A}$ في النماذج المنتهية جميعًا؛ ونعني بالمنتهي ما كانت
مجموعة عوالمه منتهية.

\begin{thm}\ollabel{thm:decidability}
  متى كان~$\Entails/ !A$ وُجد نموذج منته~$\mModel{M'}$ يحقق
  $\mSat/{M'}{!A}$.
\end{thm}

\begin{proof}
  ليكن~$\mModel{M}=\tuple{W,R,V}$ نموذجًا، ولنأخذ~$w_0\in W$ بحيث
  $\mSat/{M}{!A}[w_0]$. نسمّي~$P$ مجموعة ال!!{propositional variable}s
  الواقعة في~$!A$، ونسمي~$\Sigma$ إغلاق~$\{!A\}$ بأخذ
  الصيغ الجزئية: فـ~$\Sigma$ تضم الصيغ الجزئية
  لـ~$!A$ كلها، وهي لذلك منتهية.

  ونحدّ علاقة التكافؤ~$\sim_\Sigma$ على~$W$ كما يأتي:
  \[
    u\sim_\Sigma v
    \quad\text{إذا وفقط إذا}\quad
    \bigl(\mSat{M}{!B}[u]\ \text{إذا وفقط إذا}\
    \mSat{M}{!B}[v]\bigr)
    \quad\text{لكل }!B\in\Sigma.
  \]
  ولتكن~$[u]_\Sigma$ فئة تكافؤ~$u$. ثم نعرّف
  $\mModel{M'}=\tuple{W',R',V'}$ بهذه الحدود:
  \[
    W'=\Setabs{[u]_\Sigma}{u\in W},
  \]
  \[
    R'{[u]_\Sigma}{[v]_\Sigma}
    \quad\text{إذا وفقط إذا}\quad
    \bigl(\mSat{M}{!B}[u]\ \text{يستلزم}\
    \mSat{M}{!B}[v]\bigr)
    \quad\text{لكل }!B\in\Sigma,
  \]
  و
  \[
    V'(p)=
    \begin{cases}
      \Setabs{[u]_\Sigma}{u\in W\text{ و }\mSat{M}{p}[u]}
        & \text{إذا كان }p\in\Sigma,\\
      \varnothing & \text{خلاف ذلك.}
    \end{cases}
  \]
  وليس لاختيار ممثل الفئة أثر في هذه التعريفات. وأما~$R'$
  فانعكاسيتها وتعدّيها ظاهران. فإذا اجتمع
  $R'{[u]_\Sigma}{[v]_\Sigma}$ و
  $R'{[v]_\Sigma}{[u]_\Sigma}$، اتفق~$u$ و$v$ في صدق الصيغ كلها
  من~$\Sigma$، فيكون~$[u]_\Sigma=[v]_\Sigma$. فثبت أن~$R'$ ترتيب
  جزئي. و~$V'$ رتيبة أيضًا بالنسبة إلى~$R'$؛ فإذن
  $\mModel{M'}$ هو !!a{relational model}.

  وكل فئة في~$W'$ يعيّنها ما يصدق من صيغ~$\Sigma$ عند ممثليها؛
  ومن ذلك نحصل على الحد
  \[
    |W'|\leq 2^{|\Sigma|}.
  \]
  فـ~$W'$ منتهية على الخصوص.

  ونقيم الآن بالاستقراء البنيوي لمّة الترشيح لكل~!!{formula}~$!B$
  في~$\Sigma$:
  \[
    \mSat{M}{!B}[u]
    \quad\text{إذا وفقط إذا}\quad
    \mSat{M'}{!B}[{[u]_\Sigma}]
    \qquad(!B\in\Sigma).
  \]
  أما ال!!{propositional variable}s و$\lfalse$ و$\land$ و$\lor$
  فحالاتها تأتي بلا واسطة
  من تعريف~$V'$ وفرض الاستقراء. والمحتاج إلى بيان هو حفظ النفي والشرط
  عند هذا الاختيار لـ~$R'$.

  لنأخذ حالة~$!B\ident !C\lif !D$. إذا كان أولًا
  $\mSat{M}{!C\lif !D}[u]$، وأخذنا
  $R'{[u]_\Sigma}{[v]_\Sigma}$، فإن
  $!C\lif !D\in\Sigma$؛ وتعريف~$R'$ ينقل صدقها إلى~$v$،
  فيثبت~$\mSat{M}{!C\lif !D}[v]$. فإذا تحقق
  $\mSat{M'}{!C}[{[v]_\Sigma}]$، ثبت بالاستقراء
  $\mSat{M}{!C}[v]$؛ وانعكاسية~$R$ تعطينا حينئذ
  $\mSat{M}{!D}[v]$. ثم نطبق فرض الاستقراء ثانيةً فنحصل على
  $\mSat{M'}{!D}[{[v]_\Sigma}]$. وبهذا يثبت
  $\mSat{M'}{!C\lif !D}[{[u]_\Sigma}]$.

  وأما العكس، فلو كان~$\mSat/{M}{!C\lif !D}[u]$ لكان عالم~$v$ يحقق
  $Ruv$ و$\mSat{M}{!C}[v]$ و$\mSat/{M}{!D}[v]$. ورتابة الصدق في
  النموذج الأول تقتضي انتقال الصدق من~$u$ إلى~$v$ لكل صيغة من~$\Sigma$،
  فيثبت~$R'{[u]_\Sigma}{[v]_\Sigma}$. ويعطي فرض الاستقراء صدق~$!C$
  وعدم صدق~$!D$ عند~$[v]_\Sigma$ في~$\mModel{M'}$؛ فينتج
  $\mSat/{M'}{!C\lif !D}[{[u]_\Sigma}]$.

  وأما حالة~$!B\ident\lnot !C$ فوجهها نظير ذلك. فمن
  $\mSat{M}{\lnot !C}[u]$ و
  $R'{[u]_\Sigma}{[v]_\Sigma}$ ينقل تعريف~$R'$ صدق
  $\lnot !C\in\Sigma$ إلى~$v$. ولانعكاسية~$R$ لا يصدق~$!C$ عند~$v$؛
  فلا يصدق بالاستقراء عند~$[v]_\Sigma$ في~$\mModel{M'}$. وإن
  كان~$\mSat/{M}{\lnot !C}[u]$، وُجد~$v$ يحقق~$Ruv$ ويصدق عنده~$!C$
  أي عند~$v$. فرتابة الصدق تعطينا~$R'{[u]_\Sigma}{[v]_\Sigma}$،
  ويعطينا فرض الاستقراء شاهدًا يثبت
  $\mSat/{M'}{\lnot !C}[{[u]_\Sigma}]$.

  وقد كانت ال!!{formula}~$!A$ من~$\Sigma$؛ فتنتج لمّة الترشيح
  $\mSat/{M'}{!A}[{[w_0]_\Sigma}]$. فصار~$\mModel{M'}$ نموذجًا
  مضادًا منتهيًا لـ~$!A$، وتم المطلوب.
\end{proof}

\begin{prob}
أتمّ التفاصيل المباشرة في برهان~\olref[int][sc][dec]{thm:decidability}:
أثبت حسن تعريف~$R'$ و$V'$، وكون~$R'$ ترتيبًا جزئيًا و~$V'$
رتيبة، ثم أقم حالات المتغيرات القضوية و$\lfalse$ و$\land$ و$\lor$ في
لمّة الترشيح.
\end{prob}

وتعطي~\olref{thm:decidability} طريقًا خوارزميًا للفصل في
$\Entails !A$: نقتصر في الفحص على النماذج العلائقية التي عدد عوالمها
لا يزيد على~$2^{|\Sigma|}$، وعلى المتغيرات القضوية الواقعة
في~$!A$.

\end{document}
